\section{连续与断裂：汉帝国怎样一次次重新接线}
\label{sec:synthesis-continuity-breaks}

前二〇六年到二二〇年的汉史，很容易被切成西汉、新、东汉与三国前夜四幅彼此分开的图。每一次改朝、迁都或大战都确实改变了政治关系；但若只在断裂处停留，便看不见许多制度、道路和家庭为何能越过断裂继续运作。反过来，若只说“汉制延续”，又会抹去王莽改制、雒阳复兴、边郡军政化和区域中心兴起所造成的真实变化。本节不把连续性理解为一种不变的本质，而是追踪一组必须反复接上的连接：谁能登记土地，谁能输送粮食，谁能以皇帝名义任命官员，谁在道路中断时给予保护，谁又有资格解释秩序为何仍然有效。

\subsection{关中、长安与雒阳：首都变了，行政问题没有消失}

西汉的长安立在关中仓储、秦旧道路与开国功臣网络能够汇集的地方。首都的优势并不只在城墙和宫殿，而在于人、粮和文书可以在这里被重新分配。朝廷对关东的命令要越过函谷关、河流和众多地方关系；关东的赋役、人才与消息也要找到进入长安的路径。早汉使用封国与郡县并行的办法，正是为了让尚未完全重整的区域以不同方式接入这个中心。它不是一张事先画好的制度蓝图，而是战争之后对兵粮、名分和地方合作的暂时拼接。

长安的中心性也带来特定的脆弱。宫门附近的任命与近卫、关中仓储的调发、关东王国的兵粮，彼此一旦不能相互核验，继承或政策争论就可能迅速扩大。七国之乱之后的调整，使王国越来越难以独立动员，但没有让关东变成被动的行政末端。王国、郡县、豪强和市场继续掌握地方信息；中央只是改变了它们进入公共秩序所需经过的程序。后来的内朝与尚书台同样如此：它们使决断更快抵达文书链，也使靠近文书的人更能影响何种地方消息被称为紧急。

王莽时期的长安仍是命令与礼制的中心，却已不能仅凭旧有交通和名号保证地方响应。改币、改名、改制触及的不是抽象国家，而是市场中的交换习惯、县乡中的登记、官吏的职业路径和地方大姓的保护关系。新朝能以符命和诏书宣布新的分类，未必能使每一处地方都以同样速度接受分类。绿林、赤眉与各地武装并非从首都之外的真空中冒出；它们在不同水患、饥荒、道路和地方组织条件下，将不满、求生与政治机会连接起来。

雒阳的复兴由此不只是刘秀选择了另一座都城。它意味着新的朝廷必须让官员抵达署舍，让仓储接住粮食，让地方相信任命能够带来可预期的秩序。雒阳更接近关东的人口与交通网络，因而较容易吸纳东部的士人、商旅和官吏；它也更深地依赖这些网络。建武时重建户籍和田地登记，是让雒阳的命令重新具有重量的必要工作。若没有县乡的协作、地方中介的传递和道路的可通行性，“复汉”便只能是一套朝廷内部的称谓。

首都的转换改变了连接的重心，却没有取消连接本身。西汉的关中财政、东汉的关东士人、汉末许县对献帝的控制，都是不同集团把皇帝名义与可调度资源接在一起的方式。它们共同提醒我们，首都不是地图中央的自然点，而是人们能够在其中处理仓储、文书、任命和礼仪的一组条件。条件一旦被战争或私人网络重新分配，首都仍在，中心性却可能已经改变。

\subsection{从一纸户籍到一段迁徙：国家如何看见家庭，家庭如何避开国家}

西汉、新、东汉都需要把土地与人口写入某种可用的帐目。赋役、赈贷、兵役和移民安置无法完全脱离名册进行；但名册从来不是对家庭生活的透明摄影。一个户口数字可能包含已经外出的人、依附于大姓的人、刚刚失去劳力的人，也可能遗漏因灾荒或战争逃散的人。朝廷以户、口、田、租等类别认识社会，是为了让远方县邑可以比较和征发；家庭则在亲属、佃作、市场、婚姻和迁居之间组织生计。两种尺度经常相遇，并不总是重合。

西汉初年的休养生息在史书中常被概括为轻徭薄赋。对于经历战乱的地区，它至少意味着国家承认耕作、劳力和仓储的恢复需要时间；不能由此推断每一个家庭都获得相同的喘息。土地分散、亲属死亡、道路不通或债务存在时，减免一项征收不足以恢复原有生活。政策的意义常在于改变家庭与地方官谈判的起点：哪一类负担可以暂缓，何种损失能被承认为灾害，谁仍被认为应当留在原籍承担役使。

武帝时期的边防和财政措施，又把家庭置入更宽的物资与劳力网络。矿山的工匠、粮道的运输者、边郡的戍卒和市场中的小贩，并不生活在相同制度条件下；盐铁、钱币和转运的变化却可能通过价格、征发和采购将他们连接。把这种联系写成“国家榨取百姓”或“市场繁荣帝国”都过于平直。更准确的问题是：一项以边费为理由的政策在何地需要何种劳动，谁能通过地方关系转嫁成本，谁又因缺少关系而必须直接承担。

王莽改制使国家分类与家庭策略的冲突格外尖锐。货币、土地和官名的更易，要求地方官和市场参与者在短时间内学习新的可交换、可登记和可服从的规则。新朝的制度文本能告诉我们朝廷希望重排什么，不能逐户说明人们怎样应对。有人可能等待、投靠地方保护者、转移财物或离开原地；不同地区的道路、水患和武装环境决定这些选择并不相同。正因材料不能完整保存每一户的声音，叙事更应避免把沉默填成整齐的“民众反应”。

东汉建武度田再次让登记成为政治事件。新朝需要知道田地与人户，才能恢复赋役与赈济；被登记的一方也知道数字会影响日后的负担。地方吏员、大户与普通家庭围绕名册发生的拉扯，说明簿籍既是行政工具，也是占有关系的表面。反弹不能被简单归为地方反对秩序，因为某些反弹来自谁应承担战后缺口的争议；朝廷的统计也不能因此视为全然虚构，它记录了国家能动员何种信息与它最担心失去什么。

东汉中晚期的边地战争、灾害和党争将家庭的移动推向更广的空间。凉州、并州及其他受战争波及的地区，流民既是需要安置的人，也是地方将领、大族和州郡争取或惧怕的劳力。迁徙者带着亲族、技能、债务和旧有冲突到达新县邑；当地社会则要决定是提供土地、以依附关系收留，还是把他们视为治安与赋役问题。史书将他们常常压缩为“流民”“饥民”或“降人”，这些名称有行政用途，却不应成为他们唯一的历史身份。

汉末屯田在这一连串变化中具有双重含义。它使竞争者可以安置部分流民、组织耕作并减少军队对即时掠夺的依赖，也把人户纳入新的劳动、仓储与军政编制。不能把它描述为无条件的救济，也不宜把所有参与者只写成被迫劳动者；材料不能允许如此整齐的判断。可以确认的是，粮食、人口与军队在战乱中必须重新接合。谁能完成这种接合，谁便拥有比一纸官号更坚实的区域力量。

\subsection{礼、经与治疗：文本网络如何穿过制度断裂}

两汉的文本世界并不只有朝廷诏书。经典的抄写、讲授和校订，家族与地方的碑刻，治病与祈福的仪式，佛教相关的供养语汇、黄老与方术的实践，都让不同地区的人能够在官署之外建立可辨认的关系。它们不必然同国家对立。官方礼制、博士与太学本身就需要读书人、工匠、地方捐献者和传播者；地方仪式也可能借用帝国的年号、官名或伦理语言。将“儒家国家”与“民间宗教”写成两块彼此封闭的领域，反而看不见人们实际如何在两者之间移动。

西汉的经学制度使经典解释逐渐与官员资格相连。博士、察举与地方学校为一些人提供进入官署的路径，也使师承与评语成为跨地区流动的资源。路径不是人人都能走：旅费、家庭劳力、亲属关系和地方声望都会影响谁能求学、谁能被举荐。可是，它确实使长安的政治语言能够在地方被学习、争论和重新解释。制度的连续性，部分正靠这些携带文本和名声的人维持。

王莽时期对礼制和古典名目的高度倚重，显示文本可以为改朝提供可被理解的形式。符命、经义和官名并非无效的装饰；它们让官员知道新朝要求他们如何称呼、如何排列、如何把权力转移说成符合古义。然而，文本网络不等于执行网络。若地方官、市场与武装集团不再以同一中心为资源来源，再精致的名分也不能独自恢复道路和秩序。新朝的经验使读者看到，文本能把变革说成有先例，却不能替变革承担社会成本。

东汉的太学与石经将共同文本放进更公开的政治空间。石经残片和制度记载说明朝廷希望建立可观看、可校核的经典参照；它们不能证明所有人读同一种版本，或所有士人以相同方式理解政治。恰因经典、师承和仕进已经相连，东汉中期的名士与太学生才可能将人事不满表达成关于名教、清议和官员资格的争论。党锢不是单纯由思想引发的政治事件，也不是思想无关的派系斗争；知识网络给了地方声望进入雒阳、再从雒阳回响到地方的通道。

宗教和治疗网络的连续性更不应被简化。楚王英案中的供养语汇使某些佛教实践在东汉官方文字中可见，但不能推出全国已有统一教团；后出的高僧传记能够提供传播故事，也必须考虑其建构祖师谱系的目的。墓葬、题记、造像和传世文本各自照见局部实践。它们共同说明，求福、治病、死亡与供养的语言可以跨越州郡流动；它们不能让我们统计信众，或断言某一处仪式必然具有同一政治含义。

太平道与黄巾起事说明文本、治疗和组织关系在危机中能够转为政治动员。符号和口号给分散的人一种彼此识别、解释灾害和提出期待的语言；赋役、饥荒、地方保护缺失与军事压力，则决定这种语言在何处会凝聚为行动。所有参与者并非因同一种信仰或同一种处境而加入，也不应将所有非官方仪式倒写为叛乱的准备。更值得追问的是，为什么在某些地区，官方文书、赈济和保护已经无法提供可信的解释与帮助，而新的仪式网络可以。

汉末区域割据没有终止这些文本网络。曹氏、刘备与孙氏都需要官名、诏令、儒学语言与礼仪来任命、奖赏和说服；地方社会仍在墓葬、师承、祠祀和治疗中维持关系。二二〇年的禅代改变了谁有资格使用“汉”的最高名义，却没有让经学、礼制或宗教实践突然断绝。断裂在于，它们不再被一个中心统一收束，而要在不同政权的需要和地方社会的记忆之间继续寻找位置。

\subsection{边费如何回到家门：财政、迁徙与地方武装的反馈}

边疆政策最容易在叙事中停在边疆。白登后的和亲、武帝的北伐、河西置郡、西域经营、羌地用兵和东汉北边关系，看似发生在远离首都的地方；它们的成本却沿着谷物、布帛、钱币、劳役和人口迁徙回到内地。远征要有马、粮与道路，驻防要有传舍、仓储与守兵，互市要有可交换的物品和能被双方承认的规则。每一个环节都可能使某个县、某个市场、某个家庭的选择改变。

西汉前期的和亲与关市，显示国家有时宁愿以馈赠、贸易与守备交换一段可恢复的时间。这样的选择并非“没有成本的和平”。礼物要被筹集，使者和边郡要被维持，市场开闭会影响边民的生计与信息。可是，在关中恢复、骑兵条件有限的情况下，它也避免了把尚未稳定的资源一次耗在远征中。财政的可持续性在这里不是会计术语，而是承认一项军事选择能否让耕作、仓储和道路继续运转。

武帝朝的盐铁、均输、平准和五铢，扩大了中枢比较与调动资源的可能。它们并没有制造一座与社会隔绝的边防金库。盐矿、铁冶、转运、铸币和市场中的人，都要在政策中承担时间与风险。某处的收购可能解决军需，另一处的运输却会遇到价格、季节和地方执行的限制。《盐铁论》保存的争论之所以重要，正在于它让读者看见边费、物价、商人和民间劳力原本无法分开。

轮台诏使这种反馈变得公开。诏书承认远距征发、士卒劳苦与屯田成本有边界，并不等于西域政策立即退出历史。都护、使节、屯田和绿洲关系仍需河西道路支撑。它改变的是朝廷对可承受行动的估计。停止一项大计划，不能自动解除已经驻守者的给养、迁徙者的安置和盟友的期待；政策的后果会以更小、更难被正史完整记录的调粮、撤守和重新谈判继续出现。

赵充国对羌地的论证进一步显示，边政不是攻守二选一。山谷、水源、部众差异、粮道和季节，使屯田、招抚、分化与有限出兵各有不同作用。屯田可以减少长途运输，也需要先有安全、种子、工具和守卫；招抚能够改变关系，却不能保证所有部众作同一选择。将这些方案写成单一的“成功战略”，会掩盖它们如何把成本推向不同的人和地点。

东汉的边费反馈更加尖锐。羌地战事、北边关系和地方招募使凉州、并州及相邻地区承受持续的军事与安置压力。中央若不能稳定支付，州郡和将领就不得不向地方大族、部曲、商人和流民寻求粮、人和消息。这样的应急网络有时能使一地免于立即崩溃，却会让掌握网络的人积累军事、财政和人际资本。边费由此不只是国库支出，而成为地方权力增长的条件。

汉末的地方武装因此不能仅用宫廷失控来解释。黄巾之后的自保、凉州动乱后的招募、战争中的流民安置，都要求有人在县邑层面先行解决粮仓与治安。不同区域的道路、水网、山口和市场使这种解决办法不可能相同。冀州、关中、荆州、益州与江东各有资源与限制。区域中心的形成，正是在这些差异中把财政、迁徙和武装重新接合的过程。

\subsection{制度遗产的最后一问：谁能把后果写进共同账目}

西汉、新、东汉之间最深的连续性，不是某个官名、礼仪或税制原样存在，而是每一个政权都必须面对同一类难题：如何把地方资源送到需要它的地方，如何让命令在中途不被完全改写，如何在战争和灾害后重新登记人与地，如何使临时的保护不立刻变成私人的统治。它们的断裂也不只是国号改变，而是处理这些难题的接口在何处被截断、由谁重新占有。

王莽的改制表明，改变分类而不能取得执行网络，会使中枢的雄心反过来放大地方不确定；东汉复兴表明，恢复旧名义仍须重新修复道路、簿籍与任命；汉末的区域竞争表明，旧制度仍然有用，却已无法让所有资源服从同一套优先次序。二二〇年的受禅不是一切问题的终点。它将一个中心所能使用的最高名义转交给魏，也把许多尚未结清的户籍、边费、迁徙和地方保护问题留给不同的政权分别承担。

这也是读者面对汉史材料应保留的最后一层谨慎。帝纪和诏书往往能够告诉我们谁在何时宣布解决问题；简牍、墓葬、碑刻、宗教文本和地方叙事让某些执行的碎片可见；它们都不能替我们完成完整的共同账目。承认账目不完整，不是放弃因果解释，而是拒绝让胜者的结论、整齐的数字或一段生动传说遮住那些尚未被登记的成本。汉帝国一次次重新接线的历史，正发生在这些可见与不可见的后果之间。

\subsection{三次重建的日常：西汉、新、东汉不是同一条路的重复}

西汉初年最迫切的工作，是在秦末和楚汉战争之后让耕作、征收与任命重新有节奏。开国者面对的不是一片被彻底清空的国土，而是许多仍在运行、却彼此失去稳定联系的县邑、旧吏、地方武装和运输线。封国与郡县并存，功臣与旧有行政技术并用，正是对这种不整齐世界的应对。它让朝廷可以迅速取得协作，却把日后谁能调兵、谁能解释封国与中央关系的问题保留下来。

新朝所面对的世界已不同于开国。西汉中后期形成的货币、官名、土地关系、边防与地方声望，使任何改革都不只是恢复秩序，而是在既有秩序上重新分类。王莽的制度语言包含一种强烈的修复愿望：以古名、礼制和新名目矫正已经显露的矛盾。困难不在于人们不懂得诏书的文字，而在于旧有交易、地方保护和行政惯例不能在同一时刻被替换。改革的速度与执行网络的速度不一致，便使每一处解释权都可能成为冲突点。

东汉的重建又不能简单理解为“回到西汉”。雒阳的官署能够重新采用汉的国号、官名和经典语言，地方社会却已经历战争、迁徙与地方权力的重新排列。光武政权需要选择可以合作的豪强与地方官，需要用度田、赈济和任命处理战后人户，也需要让通往首都的道路重新可靠。复兴的成功正在于它能够把旧名义接到新的资源与人际网络上；它的长期风险也在于，网络越重要，中央越难不依赖它们的地方性利益。

三次重建并非同一循环的三个版本。西汉把开国分配转化为较稳定的中央与地方关系；新朝试图以大幅改制重排已成熟的接口；东汉则在破坏之后以复汉名义重新协商接口。它们共同证明，国家延续需要的不只是合法性的叙述，也不是只有粮仓和军队，而是让叙述、粮仓、军队和地方生活能在相近的时间上互相承认。任何一项领先太远，都会使看似完整的秩序出现裂缝。

\subsection{从边塞到都城的消息：信息不只是命令的附属}

边塞文书和地方奏报提醒我们，帝国治理不只是自上而下传递命令，也是自下而上筛选消息。居延、肩水金关与悬泉置的一片简牍可能记录某人、某马或某批粮食的出入；它的意义不在于把所有边政缩为一张表，而在于显示消息必须被写成特定格式，才可能沿着机构移动。写字的人决定怎样描述缺粮、迟到或可疑人员，上级决定何种描述值得回应。消息在每一站被压缩，最后抵达中枢的“边情”已经是多层选择的结果。

西汉朝廷经营河西和西域时，需要使节、关塞、传舍和地方官不断校正远方的信息。东汉在羌地、北边和州郡动乱中同样依赖将领与地方中介的报奏。两者的差异不只在敌人或地名，而在财政与地方组织的环境已经变化：当中枢无法稳定供给时，能报告情况的人往往也是必须自己筹粮、招兵和安置流民的人。信息与资源因此不再容易分开。报告者并非只是传声者，他可能已经成为某一片地区实际秩序的组织者。

宫廷中的外戚、宦官和近臣也应放在这个信息问题中理解。接近皇帝不只是拥有私人关系，更是能够决定哪些地方声音被迅速呈递、哪些任命附有何种解释、哪些危机被定义为可等待。内朝与尚书的存在提高了决断速度，也使决断路径变窄。政治斗争之所以容易在宫门附近加剧，正是因为这里汇集了来自不同地区、不同利益和不同时间尺度的消息。

汉末的区域中心各自建立消息与资源的回路。许县的朝廷文书、冀州与河南的粮源、荆州和江东的水路、益州的山川与地方行政，使不同集团能够较快地知道并处理本区域的问题。它们并非完全不需要外界消息，却越来越难把消息送进一套共同的优先次序。所谓帝国分裂，部分就是信息不再能与资源一起回到同一处结算。

\subsection{后果的继承：为什么二二〇年之后仍要回望汉}

汉魏禅代的仪式将一段政治关系写成有程序的终结，但许多后果并没有随之终结。屯田所处理的人口与粮食问题、边郡所留下的军政经验、太学与经学所提供的资格语言、地方大族所组织的保护关系，都要由魏、蜀、吴分别面对。后继政权使用汉代官名、文书和礼仪，并不证明它们只是汉的附属；恰恰说明这些制度接口仍是组织区域社会不可缺少的工具。

继承也意味着选择。魏在北方继承朝廷与屯田的结合，蜀地与江东则在各自的交通、人口和军事条件中选择不同的组织方式。没有一个政权可以把汉的全部资源、全部人口或全部记忆整体接收。二二〇年以后仍持续的战争、迁徙与财政压力，说明受禅解决的是最高名义的归属，不是道路、家庭和边地关系的全部问题。

因此，汉卷以二二〇年结束，并非把此后的历史排除在解释之外，而是把它作为检验：哪些连接在国号改变后仍能被使用，哪些连接已被区域化，哪些成本被后继者接手。读者在这里可以回望西汉、新、东汉的每一次重建，看到它们留下的不是单纯成功或失败，而是一组被不断转手、不断改写的可能性与负担。

\subsection{读法的回环：把跨期专题送回年序}

跨期专题的作用不是替西汉、新、东汉的年序作结论，而是让读者带着新的问题回到年序。读到七国之乱，可以追问开国的区域兵粮为何尚未完全进入共同核验；读到王莽改制，可以追问分类改变为何赶不上地方执行；读到建武度田，可以追问复兴国家怎样重新认识土地与家庭；读到羌地战争和黄巾之后的州郡动员，则要看见边费、流民与临时保护如何共同改变地方权力。每一个问题都在具体年份、地点和行动者处取得意义。

反过来，年序也限制专题的解释。财政压力不会自动产生地方割据，宗教语言不会自动产生起事，外戚或宦官的位置也不会自动使帝国失去治理能力。它们只有在特定的继承、道路、粮仓、灾害与人际选择中才构成结果。将连续与断裂放在同一卷中，是为了让读者既不把汉史写成制度自行运转的图表，也不把它写成一连串孤立的宫廷事故。

\subsection{材料在何处彼此照面}

跨期写作还要求不同材料在同一个问题上彼此照面，而不是各自替代。长安或雒阳的诏令能够告诉我们朝廷希望怎样登记、征发或安置；边塞简牍让一处传舍、关口或仓储的执行可见；墓葬、碑刻和宗教文本偶尔保留家庭、地方捐献者与仪式实践的痕迹。它们的尺度不同，恰好使读者不必把国家的语言误作全部社会，也不必把一处遗存误作全国平均。

例如，讨论流民时，帝纪的灾害与赈济记录可确定朝廷何时感到危机，地方传记可能让某些保护网络露出轮廓，屯田和军政安排则显示危机如何被转成资源问题。它们不能给出完整人数，更不能决定每个人的动机；共同的时间与空间却允许我们判断，哪些选择是在粮食、道路和安全仍有余地时作出的，哪些选择已经发生在余地极少的时刻。这样的互读，使材料限度成为解释的边界，而非叙事的终点。
